\documentclass[../../main.tex]{subfiles}
\begin{document}

\begin{flushright}
\footnotesize\textit{【自動文字起こし・要確認】}
\end{flushright}

\noindent
{\bf [解]}
$f(x)=\dfrac{x^2}{x^2+1}$，$f'(x)=\dfrac{2x}{(x^2+1)^2}$だから，$x=t$（$t\ne0$）での$C$の接線$\ell$の方向ベクトル$\vec\ell$は
\begin{align}
\vec\ell=\begin{pmatrix}1\\\dfrac{2t}{(t^2+1)^2}\end{pmatrix}
\end{align}
である．また，
\begin{align}
\overrightarrow{PQ}=\begin{pmatrix}t\\\dfrac{t^2}{t^2+1}-a\end{pmatrix}
\end{align}
だから，$PQ\perp\ell$となる時，
\begin{align}
\vec\ell\cdot\overrightarrow{PQ}=t+\dfrac{2t}{(t^2+1)^2}\left(\dfrac{t^2}{t^2+1}-a\right)=0\quad\cdots①
\end{align}

①をみたす$t\ne0$の存在条件をもとめる．$t\ne0$だから，①を変形して，
\begin{align}
a=\dfrac{1}{2}(t^2+1)^2+\dfrac{t^2}{t^2+1}
\end{align}

$y=a$と$y=(\text{右辺})$が$t\ne0$に共有点をもてばよい．$p=t^2+1$とおくと，$t\ne0\iff p>1$であり，右辺を$g(p)$とおくと，
\begin{align}
g(p)=\dfrac{1}{2}p^2+\dfrac{p-1}{p}=\dfrac{1}{2}p^2+1-\dfrac{1}{p}
\end{align}
\begin{align}
g'(p)=p+\dfrac{1}{p^2}
\end{align}

$p>0$で$g'(p)>0$だから，$g(p)$は単調増加．したがって，もとめる条件は，$p>1$の範囲で$g(p)$が$a$の値をとることと同値で，$g$が$(1,\infty)$上連続かつ単調増加であることから，
\begin{align}
a>g(1)=\dfrac{1}{2}
\end{align}

\bigskip

\end{document}
